\begin{lemma}[лемма о норме самосопряжённого оператора]\label{5.4-lemma1}
    Для самосопряжённого оператора $A \in \mathcal{L}(H)$ верно, что
    $\|A^n\| = \|A\|^n$.
\end{lemma}

\begin{proof}
    Достаточно понять, почему $\|A^2\| = \|A\|^2$. Мы знаем, что $\|A^2\| \leqslant \|A\|^2$. Распишем:
    \[
        (Ax, Ax) = (A^2 x, x) \leq \|A^2 x\| \cdot \|x\|.
    \]
    Возьмём супремум от обеих частей
    \[
        \|A^2\|_{\mathcal{L}(H)} = \sup_{\|x\|_H=1} \|A^2 x\|_H = \sup_{\|x\|=1} \|A^2 x\|_H \cdot \| x \|_H \geq \sup_{\|x\|=1} (Ax, Ax) = \sup_{\|x\|=1} \|Ax\|_H^2=\|A\|_{\mathcal{L}(H)}^2
    \]
    Получили, что $\|A^2\| \geqslant \|A\|^2$.
    В общем-то, пока мы доказали лемму только для степеней двойки (и этого достаточно для доказательства теоремы).
    Абсолютно аналогично можно показать, что лемма верна для всех $n$.
    Переход от $n$ к $n+1$ остаётся в качестве упражнения.
\end{proof}

\begin{definition}
    Функция $\angles{\cdot, \cdot}$ называется \textit{полускалярным произведением}, если для него выполняются все свойства скалярного произведения, кроме $\angles{x, x} = 0 \lra x = 0$. То есть может быть так, что $\angles{x, x} = 0$, но $x \neq 0$.
\end{definition}

\begin{statement}
    Неравенство Коши--Буняковского справедливо для полускалярного произведения (\(\left| \left\langle x, y \right\rangle \right| \leqslant \|x \| \cdot \|y \|\)).
\end{statement}

\begin{proof}
    Пусть $\angles{x, x} \neq 0$ или $\angles{y, y} \neq 0$ (пусть без ограничения общности это будет $y$). В этом случае доказательство аналогично обычному доказательству КБШ:

    Пусть $e = \frac{y}{\|y\|}$, тогда нужно доказать, что $|\angles{x, e}| \leqslant \| x \| $.

    \[
    \| x - (x, e) e\|^2 = \|x\|^2 - |(x, e)|^2 \geqslant 0
    \]

    Остаётся рассмотреть случай, когда $\angles{x, x} = \angles{y, y} = 0$.
    В вещественном случае рассмотрим выражение 
    \[
    \angles{x \pm y, x \pm y} = 0 \pm 2 \angles{x, y} + 0
    \geqslant 0
    \]

    $2\angles{x, y} \geqslant 0$ и $-2\angles{x, y} \geqslant 0 \implies \angles{x, y} = 0$.

    В комплексном случае для действительной части аналогично, а чтобы добраться до мнимой части напишем аналогичные неравенства для $x \pm iy$.

\end{proof}

\bigskip

\noindent
\textbf{Обозначения:} \(m_+=\sup\limits_{\|x\|=1} (Ax, x), \;\; m_-=\inf\limits_{\|x\|=1} (Ax, x)\)


\begin{theorem}[11.5]\label{11.5}
    Пусть $H(\C)$~--- гильбертово пространство, 
    $A$~--- самосопряжённый оператор.
    \begin{enumerate}
        \item $\sigma(A) \subset [m_-, m_+]$, 
        причём $m_-, m_+ \in \sigma(A)$, то есть
        $r(A) = \max(|m_-|, |m_+|)$
        \item $r(A)=\|A\|$
    \end{enumerate}
\end{theorem}


\begin{proof}
Не умаляя общности \(|m_+| > |m_-|\)
\begin{enumerate}
    \item
    \begin{itemize}
        \item[а)] $\sigma(A) \subset [m_-, m_+]$. 
        По теореме \nameref{11.4} спектр лежит на вещественной оси.
        Достаточно показать, что если возьмем число больше $m_+$ или меньше $m_-$, то попадем в резольвентное множество.

        Покажем, что если $\lambda > m_+ \implies \lambda \in \rho(A)$. Будем доказывать по \hyperref[11.3]{критерию принадлежности числа спектру}, то есть мы хотим 
        получить оценку вида $\|A_{\lambda} x\| \geqslant m \|x\|$.

        \[
        \|A_\lambda x \| \cdot \|x\| \geqslant |(A_{\lambda} x, x)| = |\underbrace{(A x, x)}_{\leqslant m_+ \|x\|^2} - \underbrace{\lambda (x, x)}_{> m_+ \|x\|^2}| = 
        \lambda (x, x) - (A x, x) \geqslant (\lambda - m_+) \|x\|^2
        \]

        Сокращая на $\|x\|$, получаем:
        $\|A_{\lambda} x\| \geqslant (\lambda - m_+) \|x\|$.

        \item[б)] Покажем, что $m_+$ является частью спектра. Хотим воспользоваться \hyperref[11.3]{критерием принадлежности числа спектру}\ (пункт 2): существует последовательность $x_n$, такая что $\forall n\: \|x_n\| = 1$ и $\|A_{\lambda} x_n\| \to 0$.

        Возьмём $x_n$: $(A x_n, x_n) \to m_+, \: \|x_n\| = 1$ (такие найдутся, так как $m_+$~--- супремум). 

        Можно переписать это так:
        
        \[(A x_n, x_n) - m_+ (x_n, x_n) \to 0
        \lra
        (A_{m_+} x_n, x_n) \to 0\]

        Обозначим через $B$ оператор $(m_+ I - A)$ (он ограниченный и самосопряжённый). Понятно, что $\forall x\ (Bx, x) \geqslant 0$.

        Введём полускалярное произведение: $\angles{x, y} = (Bx, y)$ и рассмотрим $\|Bx_n\|^4$:

        \[
        \|Bx_n\|^4 = (Bx_n, Bx_n)^2 = \angles{x_n, B x_n}^2 \stackrel{\text{КБШ}}{\leqslant}  
        \underbrace{\angles{x_n, x_n}}_{=(B x_n, x_n) \to 0} 
        \cdot \underbrace{\angles{B x_n, B x_n}}_{\text{ограничена}} \to 0
        \]

        $\angles{B x_n, B x_n}$ ограничена, так как

        \[
        \angles{B x_n, B x_n} = (B^2 x_n, B x_n) \leqslant \|B^2 x_n\| \cdot \|B x_n\| \leqslant \|B^2\| \cdot \|x_n\| \cdot \|B\| \cdot \|x_n\| = \|B\|^3
        \]
    \end{itemize}
    
    \item По теореме \nameref{10.1} $r(A) = \lim\limits_{n \to \infty} \sqrt[n]{\|A^n\|}$. Применяя  \hyperref[5.4-lemma1]{лемму о норме самомопряжённого оператора}, получаем требуемое.
  \end{enumerate}
  
\end{proof}